\section{Introduction}
\label{sec:introduction}

Grid load forecasts and energy demand forecasts have inspired a growing body of statistics and machine learning literature, in the interest of Transmission System Operators (TSOs), international organizations, regulatory bodies, and other related entities. In the European context, the European Network of Transmission System Operators for Electricity (ENTSO-E) heavily depends on mid-term demand predictions to conduct mid-term adequacy forecasts \cite{midterm_adequacy} and seasonal outlook overviews \cite{nseasonal}. Forecasts of up to half a year are essential here for coordination at the national, regional, and European levels \cite{entsoe_mop}. This coordination supports risk preparedness across all three levels, facilitating the refinement of dedicated Risk Preparedness Plans involving governments and national regulatory authorities, and the maximization of cross-border capacities \cite{entsoe_mop}. To this end, the ENTSO-E systematically collects load forecasts from TSOs, including day-ahead, week-ahead, month-ahead, and year-ahead predictions, on their Transparency Platform \cite{entsoe_mop}. 

The necessity of high-quality medium-term forecasts, alongside daily and long-horizon predictions, is essential to ensure security of energy supply, and to support market stability \cite{curiel_integrating_2025}. Improved medium-term forecast accuracy supports renewable energy source integration, resource adequacy assessments and resilient grid operations, particularly in the current scenario of rising grid congestion, increased dependence on renewable energy, and heightened climate and economic variability \cite{curiel_integrating_2025}. Predictions at mid-term horizons can further help protect vulnerable populations from shortages and price volatility \cite{van_de_sande_enhancing_2024}.

At the Dutch national level, TenneT, the sole TSO in the Netherlands, expresses interest in short-, mid-, and long-horizon forecasts but strictly enforces accuracy constraints, noting that ``grid operators consider forecasts with maximum deviations of 5\% to 10\% to be usable and accurate'' \cite{tennet_forecasts}.

Current grid load forecasting models for the Netherlands rely on historical load data, standard meteorological data \cite{curiel_integrating_2025,van_de_sande_enhancing_2024,ashtar_hybrid_2025}, and economic indicators \cite{curiel_integrating_2025,van_de_sande_enhancing_2024,ashtar_hybrid_2025}, failing to capture the high-dimensional socio-economic signals that as per recent research \cite{chen_monthly_2025,r_forecasting_2025,borunda_machine_2025} may be encoded in VIIRS \textit{Nighttime Light} (NTL) imagery collected via satellite \cite{nasa_viirs_vnp46a2}. This study presents experiments evaluating the inclusion of this data in grid system load forecasting for the Netherlands, with the aim of investigating forecasting methods potentially more suitable to the needs of TenneT, ENTSO-E and other TSOs.

Within the fields of energy informatics, sustainability, and remote sensing, research has emerged that validates the potential of \textit{Multi-Modal Learning} (fusing heterogeneous data such as text, images, and time-series), specifically incorporating VIIRS nighttime light imagery, in improving grid load forecast accuracy \cite{chen_monthly_2025,r_forecasting_2025,borunda_machine_2025}. Furthermore, attention-based mechanisms have been tested as means of capturing complex intra- and inter-modal relationships to enhance time series forecasting for other domains, such as energy generation forecasting \cite{schubnel_solarcrossformer_2025} or supply chain demand forecasting \cite{wang_causal-aware_2025}.

This study presents the design and evaluation of an intra- and inter-modal attention-based mechanism, adapting architectures such as \cite{schubnel_solarcrossformer_2025,wang_causal-aware_2025}, that fuses NTL and time-series modalities for the Dutch context. This proposed model is compared against established baseline models previously tested specifically for prediction of grid system load in the Netherlands --- Prophet-LSTM \cite{van_de_sande_enhancing_2024}, Seq2Seq \cite{ashtar_hybrid_2025}, and Multiple Linear Regression --- complemented by na\"ive reference baselines (24-hour persistence, seasonal-na\"ive, and drift). Furthermore, the study presents step-by-step inclusion/exclusion studies, as done in reference literature such as \cite{wang_causal-aware_2025}, to quantify the specific predictive gain attributed to the inclusion of VIIRS imagery and attention layers. Existing multi-modal architectures from other problem domains, such as SolarCrossFormer \cite{schubnel_solarcrossformer_2025} and Causal-Aware Multimodal Transformer \cite{wang_causal-aware_2025}, as well as their associated experiment designs, serve as reference material adapted here for the Dutch grid load forecasting problem statement. Accessible and well-validated public datasets from the ENTSO-E \cite{entsoe_power_stats}, the U.S. National Aeronautics and Space Administration (NASA) Black Marble project \cite{nasa_viirs_vnp46a2}, the Dutch Royal Netherlands Meteorological Institute (KNMI) \cite{knmi_hourly_obs}, and the Dutch Central Bureau of Statistics (CBS) \cite{cbs_statline}, supplemented by algorithmically derived Dutch holiday calendars \cite{holidays_python,openholidays_api}, serve as the basis for the work; no novel data collection is presented. This should allow for ease of future reconstruction and validation of the given experiments.

\bigskip

\noindent\textbf{Research Question:} \textit{How does the performance of a multi-modal attention-based deep learning model, integrating VIIRS nighttime light data, economic indicators, and weather data, compare to current unimodal (time-series only) baselines for electricity grid load forecasting in the Netherlands?}

The following sub-questions are formulated:
\begin{itemize}
    \item \textbf{SQ1 (Architectural Adaptation):} How can a multi-modal attention-based architecture be designed to fuse information from VIIRS satellite imagery with economic and weather time-series data for the specific context of the Dutch power grid system load prediction?
    \item \textbf{SQ2 (Benchmarking):} How does the proposed multi-modal model compare in performance to established baselines for the Netherlands (specifically Prophet-LSTM, MLR, Seq2Seq) in terms of predictive accuracy over 60 day to 180 day horizons?
    \item \textbf{SQ3 (Inclusion/Exclusion Analysis):} What is the specific predictive contribution of 2D VIIRS nighttime light imagery and cross-modal attention mechanisms in system grid load prediction, as opposed to unimodal self-attention mechanisms with or without 1D aggregate NTL data?
\end{itemize}
